\documentclass[12pt]{article}

% ---------- Packages ----------
\usepackage{geometry} 

\usepackage[utf8]{inputenc}
\usepackage[T1]{fontenc}
\usepackage{lmodern}
\usepackage{amsmath, amssymb, amsthm}
\usepackage{enumitem}
\usepackage{fancyhdr}
\usepackage{titlesec}
\usepackage{hyperref}
\usepackage{xcolor}
\usepackage{graphicx}
\usepackage[breakable]{tcolorbox}
\usepackage{booktabs}
\usepackage{algorithm}
\usepackage{algpseudocode}


% ---------- Page Setup ----------
\geometry{margin=0.8in}
\setlength{\parskip}{0.3em}
\setlength{\parindent}{0pt}

% ---------- Header & Footer ----------
\pagestyle{fancy}
\fancyhf{}
\rhead{Intro to Algorithms Notes}
\lhead{Your Name}
\cfoot{\thepage}

% ---------- Section Formatting ----------
\titleformat{\section}{\Large\bfseries}{\thesection}{1em}{}
\titleformat{\subsection}{\large\bfseries}{\thesubsection}{1em}{}
\titleformat{\subsubsection}{\normalsize\bfseries}{\thesubsubsection}{1em}{}

% ---------- Custom Environments ----------
\newtheorem{definition}{Definition}[section]
\newtheorem{theorem}{Theorem}[section]
\newtheorem{lemma}{Lemma}[section]
\newtheorem{corollary}{Corollary}[section]
\newtheorem{proposition}{Proposition}[section]

\tcbset{exampleboxstyle/.style={ % Define a style for tcolorbox
  colback=blue!5!white,
  colframe=blue!75!black,
  title=Example,
  breakable % The breakable key is here
}}

\newenvironment{examplebox}[1][]{ % Now define a standard newenvironment
  \begin{tcolorbox}[exampleboxstyle,#1] % Apply the style and allow extra options
}{
  \end{tcolorbox}
}

% ---------- Title ----------
\title{\textbf{Intro to Algorithms Notes}\\[1ex]\large Based on Cormen et al.}
\author{Abdullah Yassine}
\date{\today}

\begin{document}

\maketitle
\tableofcontents
\newpage

% ---------- Chapter Example ----------
\section{Chapter 1: The Role of Algorithms in Computing}

\subsection{What is an Algorithm?}
\begin{definition}
An \textbf{algorithm} is a freaky well-defined computational procedure that takes some value or set of values as input and produces some value or set of values as output.
\end{definition}

\subsection{Algorithmic Efficiency}
Discuss time complexity and growth of functions. which is great for thinking

\subsubsection{Asymptotic Notation}
\begin{itemize}
  \item Big-O: $f(n) = O(g(n))$ means ...
  \item $\Omega$, $\Theta$ also covered.
\end{itemize}

\subsection{Example Algorithm}
\begin{algorithm}[H]
\caption{Insertion Sort}
\begin{algorithmic}
\For{$j = 2$ to $A.length$}
  \State $key = A[j]$
  \State $i = j - 1$
  \While{$i > 0$ and $A[i] > key$}
    \State $A[i + 1] = A[i]$
    \State $i = i - 1$
  \EndWhile
  \State $A[i + 1] = key$
\EndFor
\end{algorithmic}
\end{algorithm}

\begin{examplebox}
Given the input array $[5, 2, 4, 6, 1, 3]$, Insertion Sort produces ... as
\end{examplebox}

\subsection{Key Takeaways}
\begin{itemize}
  \item Algorithms are not just code—they are general, logical procedures. again
  \item Efficiency is crucial: compare brute-force vs. optimized versions. never do this again
\end{itemize}

\newpage
\section{Chapter 2: Getting Started}
% ... continue structure similarly

\end{document}
